\documentclass[11pt]{article}
\usepackage[T1]{fontenc}
\usepackage{lmodern,amsmath,amssymb,amsthm,geometry,microtype,xcolor}
\usepackage[colorlinks=true,urlcolor=blue!55!black,linkcolor=black]{hyperref}
\geometry{margin=.93in}
\setlength{\emergencystretch}{2em}
\newtheorem{theorem}{Proposed theorem}
\newcommand{\dd}{\,dl}
\newcommand{\dt}{\,dt}
\newcommand{\Rr}{\mathcal R}
\newcommand{\Ll}{\mathcal L}
\newcommand{\E}{\mathbb E}
\newcommand{\eps}{\varepsilon}
\newcommand{\norm}[1]{\lVert #1\rVert}
\newcommand{\repo}{https://github.com/TodManlaibaatar/OOD-reversals-theory/blob/12aa3b6f48f83a097600bcf9b44d93399e2f9697}
\newcommand{\src}[2]{\href{\repo/09_canonical_regime_AN04/AN04_#1_v1.tex}{#2}}
\title{AN04: From the first small weak response to the central E5 depth comparison}
\author{Version 1 \quad|\quad Research draft for review}
\date{5 November 2026}
\begin{document}
\maketitle
\begin{abstract}
At the central critical initialization, a fixed small weak-response hit has bounded total mass.
Dissipation there gives a fixed loss gap and places the first measured half response within bounded time.
We present the argument for carrying the signed strong-minus-weak history of exposed fourth-quadrant labels to the coherent witness time. Its quantitative target is a negative vertical probe output of larger order than both the horizontal cost and all adverse output. The conclusion concerns one endpoint of the initialized flow. The companion review identifies the estimates on which the proposed conclusion depends.
\end{abstract}
\section{Claim, clocks and source results}\label{an04ce5v1:setup}
Use normalized time and the equally weighted mixture of $P_1=N(e_1,q^2I)$ and $P_2=N(\rho e_2,q^2I)$ with identity target $D$. Write $\E_j$ for expectation on $P_j$ and $\E$ for mixture expectation.
Fix $\rho=2/3$, $\zeta=5\pi/9$, $\lambda=\rho^2$, $L=\lambda/2$, $\omega=(1-\lambda)/2$ and $\ell=\log(1/q)$. With $c>0$ fixed take $S_0=cq^{18/5}$. The initial vectors are $U_\alpha=W_\alpha=\sqrt{S_0}(\cos\alpha,\sin\alpha)$, with $dl=d\alpha/(2\pi)$. The quadrant and half-circle names always refer to this initial label. With $a_+=\max(a,0)$, put
\[
 \begin{gathered} f(D)=\int W(U^\top D)_+\dd,\qquad r=D-f,\qquad
 S=\int |U|^2\dd=\int |W|^2\dd,\\
 \Ll_j=\tfrac12\E_j|r|^2,\quad \Ll=(\Ll_1+\Ll_2)/2.\end{gathered}
\]
Balance holds label by label. For $\Theta=(U,W)$ the exact characteristics are
$U'=\mathsf R^\top W$, $W'=\mathsf R U$ with
$\mathsf R=\E[rD^\top{\bf1}_{U^\top D>0}]$ for the receiving label.
Thus $\Ll'=-\|\Theta'\|^2$ and $S'=2\E[f\cdot r]\le \E|D|^2/2$.
All unlabeled norms of parameters are in $L^2(dl)$. Set $Q_1=1+q^2$, $Q_2=Q=\lambda+q^2$ and
\[
 b_j=\E_j[D_j(U^\top D)_+]/Q_j,\quad v_j=\nabla_U b_j,\quad
 \theta_j=(v_j)_j,\qquad
 F=\begin{pmatrix}m&h\\z&n\end{pmatrix},\quad Fe_j=\int W b_j\dd.
\]
Let $t_h$ be the first half mass of right ancestry and $t_a=\inf\{t\ge t_h:n(t)=a\}$ for a weak response level $a$. Thus $t_{1/2}$ is measured from the actual flow. The central coherent witness has offset $s_*$, comparison interval $J$, and $e_*=\min_{s\in J}\mathcal E_{\rm coh}(s)=\mathcal E_{\rm coh}(s_*)<0$. Define $\xi=(\sin(q\zeta),-\cos(q\zeta))$ and $T=t_{1/2}+s_*$.
\begin{theorem}[Central initialized endpoint]\label{an04ce5v1:main}
For every fixed $c>0$ and all sufficiently small $q$, the first half response is finite. At $T$,
\[
 m\ge\tfrac12,\qquad n\ge\tfrac14,\qquad
 \Ll_j(T)/\Ll_j(0)\le169/196\quad(j=1,2).
\]
The proposed depth estimate is
$e_q(T)=(|\xi-f_T(\xi)|^2-1)/(2q^2)\to-\infty$ as $q\downarrow0$.
In particular $e_q(T)<e_*$. The unscaled probe loss tends to $1/2$.
\end{theorem}
\paragraph{Inputs from the theory repository.}
\src{trained_half_mass_bridge_donors}{HBD} and \src{bridge_return_before_weak_learning}{BRW}
give $t_h=\log(4/S_0)+o(1)$. At $t_h+b$ for each fixed $b\ge0$, the strong mass and regression tend to $M(b)=(1+e^{-b})^{-1}$, the right radial density tends to $4M(b)\cos^2\alpha$, and $4\Ll\to (1-M(b))^2+\lambda$. The Cartesian entrance envelopes have horizontal scale $H_0=\sqrt{S_0}e^{(1/2+q^2)t}$ and vertical scale $V_0=\sqrt{S_0}e^{(L+q^2)t}$. In particular both second components are $O_b(q)$ and the left radial mass per label is $O_b(q^2(1+t_h)^2+S_0)$ at fixed delay.
\src{core_output_descent}{COD} gives the resident rank ceiling and bottom-path bounds.
\src{transverse_balance_and_weak_clock}{TBW} gives the seed
\[
 q^{-2}\int_{\rm Q2}[(U_2+W_2)/2]^2\dd\ \longrightarrow\ c^{2\omega}(4e^b)^\lambda/8 .
\]
On compact Q2 interiors the seed divided by $q$ tends to $c^\omega(4e^b)^L\sin\alpha$ and both first components are $o_b(q)$. Fix $\eta$ as in TBW's small-hit proof, decreasing it if needed. On a fixed positive-length compact Q2 arc before $t_\eta$, $v=(U_2+W_2)/2>0$, $|U_1|,|W_1|,|U_2-W_2|/2\le v/2$, and $v'\ge\lambda v/4$ after a sufficiently late fixed entrance. Its loss budget gives $m>99/100$ there.
\src{residual_exchange_channels}{REC} and \src{cap_free_cross_response}{CFCR} give the Gaussian regression flow and shape estimates. We use the transverse energy identity of \src{gate_transport_and_transverse_memory}{GTM}. The \href{\repo/04_witnesses_persistence/AN06_learned_orbit_witnesses_v2.tex}{WIT central margins} give $s_*>-\log(3)/\lambda$. Constants may depend on $c$, on fixed witness data and on choices already fixed. The following passage and signed estimates are the work submitted for review.

\section{Mass and time at the first small hit}\label{an04ce5v1:mass}
Write $t_0=t_h+b_0$ for a sufficiently late fixed entrance, $s=t-t_0$,
$\tau=t_\eta$, and $E=\tfrac12\int (U_2^2+W_2^2)\dd$. Stop at the first of
$\tau$, $S=4$, and $t_d=t_0+(2/Q)\log(d/q)$ for fixed $0<d<1$.
On this interval
\begin{equation}\label{an04ce5v1:energy}
 E'\le QE+C(q\sqrt E+q^2),\qquad |n|\le E+Cq\sqrt E .
\end{equation}
Indeed $b_2=(U_2)_++O(q|U|)$ gives the second bound by Cauchy--Schwarz.
The exact transverse identity is $E'=\E[D_2r^\top\partial_2 f+r_2 f_2]$.
On $P_1$, $\|\partial_2 f\|_\infty\le\sqrt{2SE}$, $D_2$ has norm $q$, and $r_2 f_2\le D_2^2/4$. On $P_2$, Euler's identity for the homogeneous map $f$ writes the integrand as
\[
 (D_1-f_1)(f_1-D_1\partial_1 f_1)+2r_2 f_2-D_1r_2\partial_1 f_2.
\]
The first product is at most $D_1^2(1-\partial_1 f_1)^2/4$ by completing the square. The last term costs $O(q\sqrt E)$, and the middle term has expectation at most $2Qn$. The cluster weight $1/2$ gives \eqref{an04ce5v1:energy}. Dividing by $2\sqrt{E+q^2}$ and integrating yields
\begin{equation}\label{an04ce5v1:envelope}
 \sqrt E\le C_{b_0}q e^{Qs/2},\quad
 \int_{t_0}^t\sqrt E\le C_{b_0}d,\quad \int_{t_0}^t E\le C_{b_0}d^2 .
\end{equation}
Projecting $S'$ on the principal cluster coordinates and completing the two residual squares gives $S'\le Q_1m(1-m)+Qn(1-n)+q^2/2$. Moreover
\[
 m'=a_m(1-m)+O(\sqrt E+q),\qquad
 a_m=\tfrac{Q_1}2\|\nabla m\|^2\ge Q_1m^2/S.
\]
Here $|\langle\nabla m,\nabla n\rangle|\le Cq S$ since $|v_1^\top v_2|\le Cq$, $|h|+|z|\le C(\sqrt E+q)$, and the projected shape error is $O(q)$ on the stop. Euler's identity $\langle\Theta,\nabla m\rangle=2m$ gives the lower bound. Integrating the damped equation bounds the positive deficit cost by $C(1-m(t_0))_++C\int(\sqrt E+q)$. Hence
\[
 S(t)\le S(t_0)+C(1-m(t_0))_++C_{b_0}(d+d^2)+o_q(1)<4 .
\]
Choose $b_0$ first using $S(t_0),m(t_0)\to M(b_0)$, then $d$, then $q$. This removes the mass stop.

On the TBW cone the weak gate misses an exponentially small data moment, so $\mathsf R_{22}\ge Q(1-n)/2-Cq$. The off-diagonal entries cost $C(\sqrt E+q)$. Thus $v'\ge[Q/2-C(E+\sqrt E+q)]v$. If the hit is later than $t_d$, integration of \eqref{an04ce5v1:envelope} gives $v(t_d)\ge k_{b_0,d}>0$ on the arc. Beyond $t_d$, TBW's cone and the global linear mass bound imply $|\mathrm{arc}|\, k_{b_0,d}^2 e^{\lambda u/2}\le4+Cu$ until the hit, where $|\mathrm{arc}|$ denotes label measure. This forces $u=O(1)$. If the hit is earlier, $\eta\le E(\tau)+Cq\sqrt{E(\tau)}$ gives the opposite clock inequality from \eqref{an04ce5v1:envelope}. In either case
\begin{equation}\label{an04ce5v1:smallhit}
 \tau-t_h=\ell/L+O(1),\qquad
 \sup_{[t_0,\tau]}S+\int_{t_0}^{\tau}(|1-m|+|n|)\dt=O(1).
\end{equation}
For the absolute deficit use both signs in the damped equation. On an additional fixed interval the linear mass bound and bounded regressions suffice.

\section{A fixed loss gap and the first half response}\label{an04ce5v1:half}
Let $A=\operatorname{diag}(\sqrt{Q_1},\sqrt Q)$, $u=A(b_1,b_2)^\top$, and $D_{\Ll}=4\Ll$.
Cluster projection gives $D_{\Ll}=\|(I-F)A\|_F^2+\chi$ with $0\le\chi\le Cq^2(1+S)^2$. Indeed $r$ is $(1+S)$-Lipschitz on each cluster, and approximation by its mean-center value times $D_j/\E_jD_j$ leaves $O(q(1+S))$ in $L^2(P_j)$. Since $FA=\int Wu^\top\dd$, $\int uu^\top\dd\succeq (FA)^\top(FA)/S$. The output gradient is $\dot W=(I-F)Au/2+N_W$ with $\|N_W\|\le Cq\sqrt S(1+S)$. Consequently
\begin{equation}\label{an04ce5v1:gram}
 \|\dot W\|\ge\frac{\sigma_{\min}(FA)}{2\sqrt S}\|(I-F)A\|_F-Cq\sqrt S(1+S).
\end{equation}
The bounded mass in \eqref{an04ce5v1:smallhit} bounds $|n'|$ uniformly. There is a fixed $a>0$ with $\eta/2\le n\le\eta$ on $[\tau-a,\tau]$. For every sufficiently late fixed $b$ the entrance at $t_h+b$ precedes this interval. Loss monotonicity and then $b\to\infty$ give $\limsup_{q\to0}(D_{\Ll}(\tau-a)-Q)\le0$. With $\gamma=\sqrt{Q/Q_1}$, the exact square completion reads
\[
 2\gamma\det F=(Q-D_{\Ll}+\chi)/Q_1+
 (m-1+\gamma n)^2+(z-\gamma h)^2+2\gamma(1-\gamma)n .
\]
Thus $\det F$ is bounded below on a fixed time interval. Since $\|FA\|$ is bounded there, its least singular value is bounded below too. The error norm in \eqref{an04ce5v1:gram} is at least $\sqrt Q(1-\eta)$, so $D_{\Ll}'\le-c_0<0$ there. In particular
\begin{equation}\label{an04ce5v1:gap}
 D_{\Ll}(\tau)\le Q-\Delta
\end{equation}
for a fixed $\Delta>0$. Afterwards $\sigma_{\min}(FA)\ge\sqrt Q-\sqrt{Q-\Delta}$ by the triangle inequality. Before half response the remaining error norm is at least $\sqrt Q/2$. On $[\tau,\tau+K]$, $S\le B+C(t-\tau)$. Choose $K$ so that the dissipation lower bound proportional to $\int_0^K (B+Cu)^{-1}du$ exceeds all available loss, then take $q$ small enough to absorb the error in \eqref{an04ce5v1:gram} throughout that fixed interval. This forces
\begin{equation}\label{an04ce5v1:halfclock}
 0\le t_{1/2}-\tau\le K,\qquad t_{1/2}=t_h+\ell/L+O(1).
\end{equation}
The argument requires neither monotonicity of $n$ nor a small cross-response cap.

\section{Axis separation and the learning guards}\label{an04ce5v1:guards}
Fix a compact arc $J_2\Subset{\rm Q2}$ and a bounded half-centered window with right endpoint $t_f=t_{1/2}+K$. Charge vertical components outside $J_2$ and horizontal components inside it. With $P$ the corresponding orthogonal projection on the two-layer parameter space, put $E_J=\|P\Theta\|^2/2$. On the bounded-mass passage,
\begin{equation}\label{an04ce5v1:defect}
 E_J'\le[Q+C(|n|+|1-m|)]E_J+C(q\sqrt{E_J}+q^2).
\end{equation}
To see the coefficient, pair the four-regression flow with $P\Theta$. For the $n$-gradient the pairing is $2\int_{J_2^c}W_2(U_2)_+\dd+O(q\sqrt{E_J})$ using $b_2=(U_2)_++O(q|U|)$, $|U_2(\theta_2-{\bf1}_{U_2>0})|\le Cq|U|$, and $|(v_2)_1|\le Cq$. Its absolute value is bounded by $2E_J+Cq\sqrt{E_J}$. Interchanging coordinates gives the corresponding bound for $m$. For $z,h$ the pairing is respectively $z+O(qS), h+O(qS)$: one charged coordinate differentiates the output or its matching feature on each side of $J_2$. The leading cross terms are negative squares, and $|z|+|h|\le C(\sqrt{E_J}+q)$. Pairing the shape remainder costs $Cq\sqrt{E_J}$. This gives \eqref{an04ce5v1:defect} without a coefficient depending on the arc.

Apply the square-root comparison to $E_J+q^2$. By \eqref{an04ce5v1:smallhit} and \eqref{an04ce5v1:halfclock}, for each fixed $b\ge b_0$,
\[
 \sup_{t_h+b\le t\le t_f}\sqrt{E_J(t)}
 \le C_K e^{-Lb}\{1+\sqrt{E_J(t_h+b)}/q\}
\]
for small $q$ (enlarging $C_K$ to absorb $q^2\ell$). At fixed $b$, the right vertical energy divided by $q^2$ tends along extraction to $\tfrac12\int (x^2+y^2)d\mu$ for the BRW resident chart. Indeed this holds on compact interiors by tagged transfer, and the HBD Cartesian bound $O_b(q)$ permits exhaustion at each fixed $b$. Its product with $e^{-2Lb}$ tends to zero as $b\to\infty$: for positive parts COD bounds the discounted squares at right rank $f$ by $C(1+f^{-1/3})^2e^{-\lambda f b/8}$, integrable in $df$ on Q1, with $f=1/2$ on Q4. Order and the bottom path bound control $x_-$ by $C\sqrt{1+b}$. From the bottom $y$ equation and the BRW mean budget, $y_-\le C(1+\int_0^b e^{Lu}(1+u)^{-2/3}du)=o(e^{Lb})$.
At fixed $b$ a compact Q3 interior has radial mass $O_b(S_0)=o(q^2)$ (COD's trained transfer); its endpoint strips may again be discarded by the Cartesian bound. On $J_2$ the charged horizontal components are $o_b(q)$. On the rest of Q2, compact transfer, balance and the seed density give a discounted squared cost bounded by $C\int_{{\rm Q2}\setminus J_2}\sin^2\alpha\dd$. Thus one first takes $q\downarrow0$ at fixed arcs and delay, then $b\to\infty$, then exhausts Q2. No fixed-delay chart is evaluated at delay $\ell/L$.

In particular the right vertical norm, $|z|+|h|$, and $|T_1|$ with $T_1=\int W_1W_2|v_1|^2\dd$ vanish uniformly on $[t_0,t_f]$, combining the late estimate with the entrance bound up to any fixed delay. Away from the entrance, $m$ has positive damping: up to $\tau$ use TBW, and after $\tau$ use $Q_1(1-m)^2\le Q-\Delta$. The scalar $m$ equation with vanishing cross forcing yields $m\to1$ uniformly on compact half-centered windows.

For the weak gradient the exact identity is
$\|\nabla n\|^2=\int(b_2^2+W_2^2|v_2|^2)\dd$.
On each compact $J_2$, before the small hit the positive cone and balance give $|W_2-U_2|^2\le|W_1^2-U_1^2|$. On a bounded continuation the difference solves a scalar equation with bounded coefficient $-\mathsf R_{22}$ and an off-diagonal forcing bounded in norm by a constant times the horizontal norm. Positivity up to an $L^2$ error follows by the same comparison for $(U_2+W_2)/2$ with coefficient $\mathsf R_{22}$. Thus the norm errors can be bounded by the charged defect and taken to zero in the exhaustion above. The cone is applied on each fixed compact separately, using its strict TBW entrance. Outside $J_2$ the charged vertical energy bounds the error. Since $|b_2(1-|v_2|^2)|\le Cq|U|$, this gives $\|\nabla n\|^2=2n+o(1)$ on each compact half-centered window. The exact scalar regression equation there reduces to $n'=Qn(1-n)+o(1)$. Using the measured centering $n(t_{1/2})=1/2$ gives uniformly for bounded $s$
\begin{equation}\label{an04ce5v1:logistic}
 (m,z,h,n)(t_{1/2}+s)\longrightarrow(1,0,0,N(s)),\qquad
 N(s)=(1+e^{-\lambda s})^{-1}.
\end{equation}
This also bounds $1-n$ away from zero through $T$ by combining the fixed continuation and the pre-hit bound. At $s_*$, $N(s_*)>1/4$. Cluster projection and $f_0(D)=S_0D/4$ give loss ratios tending to $0$ and $(1-N(s_*))^2<9/16<169/196$. This proves both response and both loss guards.

\section{Probe exposure and the signed history to be estimated}\label{an04ce5v1:history}
Set $p_\alpha=(U_\alpha(T)^\top\xi)_+$, $F_i=f_{T,i}(\xi)$, $H=F_1/q$ and $w=Q(1-n)$. Before $t_d$ on the stopped interval, \eqref{an04ce5v1:envelope} bounds $\int\sqrt E\dt$. After that interval ends there are at most $O(1)$ units left to $T$, by the clock comparison. Consequently
\[
 \int_{t_0}^T(\|\Theta'\|+\|r\|_{P_1})\dt\le C,\qquad
 \|\mathcal U\|\le C,\quad
 \mathcal U_\alpha=\sup_{[t_0,T]}(|U_\alpha|+|W_\alpha|).
\]
Here the regression terms in $\|\Theta'\|$ are bounded by $C(|1-m|+|h|+|z|+\sqrt E+q)$ since $\|\nabla n\|\le C(\sqrt E+q)$. The strong residual costs $C(|1-m|+|z|+q)$ by shape projection. Minkowski's inequality bounds the time supremum in label norm.
On an excursion below $U_2=-A_0q\sqrt\ell\,\mathcal U_\alpha$, the $P_2$ gate has Gaussian upper tail and costs at most $Cq^2\mathcal U_\alpha$ per time in $U_2'$, for a sufficiently large fixed $A_0$. The $P_1$ cost is at most $Cq\mathcal U_\alpha\|r\|_{P_1}$ by Cauchy--Schwarz. Integrate from the last crossing or $t_0$, using the $O_{b_0}(q)$ entrance. Uniformly up to $T$,
\[
 (U_2)_-\le C_{b_0}q+Cq\sqrt\ell\,\mathcal U_\alpha,\qquad
 \|p\|\le Cq\sqrt\ell,\qquad |F_1|+|F_2|\le Cq\sqrt\ell .
\]
Define the strong-minus-weak history and the signed budget by
\begin{align}
 G_\alpha&=\int_0^T (Q_1 z b_1-w b_2)\dt,\nonumber\\
 M_q&=\frac{\int_{\rm Q4}pG_\alpha\dd-\int_{{\rm Q4}^c}p(-G_\alpha)_+\dd}{2q^2}.
 \label{an04ce5v1:budget}
\end{align}
The sign follows from $W_2'=(w b_2-Q_1zb_1)/2+\varepsilon_\alpha$. Orthogonal projection of both the residual and feature off $D_j$ in each cluster bounds $|\varepsilon_\alpha|\le Cq^2(1+S)|U_\alpha|$. Bounded mass through $T=O(\ell)$ therefore gives
$\int p|W_2(T)+G_\alpha/2|\dd\le C\|p\|(\sqrt{S_0}+q^2\ell)=o(q^2)$.
In particular $-F_2/q^2\ge M_q-o(1)$. The exact scaled loss is
\[
 e_q=H^2/2-\frac{\sin(q\zeta)}q H+
 \left(\cos(q\zeta)+F_2/2\right)F_2/q^2 .
\]
The session's sufficient criterion (3) is
\begin{equation}\label{an04ce5v1:criterion}
 \liminf_{q\downarrow0}\left[M_q-H^2/2+\frac{\sin(q\zeta)}qH+e_*\right]>0 .
\end{equation}
For this implication, $A_q=H^2/2-q^{-1}\sin(q\zeta)H-e_*=O(1+\ell)$ and the vertical coefficient is $1+O(q\sqrt\ell)$. On the branch $F_2=o(1)$ the threshold for $-F_2/q^2$ is $A_q+o(1)$, since $q\sqrt\ell(1+\ell)=o(1)$. The estimates below will give a diverging margin in \eqref{an04ce5v1:criterion}.

\paragraph{Why the signed history is needed.}
There is also an exact pairwise diagnostic:
\[
 zF_1-mF_2=-\tfrac12\iint\det(W_a,W_b)
   (b_{1,a}p_b-b_{1,b}p_a)\,dl_a\,dl_b .
\]
It follows by antisymmetrizing the two products of integrals. To see the quantitative issue, prescribe two equal strong masses with scaled input angles $\zeta-\epsilon,\zeta+\epsilon$ and output angles $-1,+1$, respectively, and a vertical weak population. In the limiting probe formula only the first strong atom is exposed. The horizontal output divided by $q$ is $\epsilon/2$ and the vertical output divided by $q^2$ is $-\epsilon/2$. The scaled error is $\epsilon^2/8-(\zeta+1)\epsilon/2$, which approaches zero as $\epsilon\downarrow0$. A favorable pairwise sign can therefore coexist with an error above $e_*$. The initialized proof needs \eqref{an04ce5v1:criterion} with its magnitude and its charge for adverse labels.

For the remainder set $A_R=\{\cos\alpha>0\}$, $A_L=\{\cos\alpha<0\}$,
\[
 C(t)=\int (U_1)_+(U_2)_+\dd,\quad I(t)=\int_{t_0}^t C(u)\,du,\quad
 B(t)=q(1+t-t_0),\qquad t_R=t_h+\ell/L-R .
\]
Here $C(t)$ denotes the overlap and unindexed constants $C$ in inequalities denote bounds. For all sufficiently large fixed $R$, $t_R$ precedes both $t_d$ and $\tau$ by the two-sided clock estimate. Equation~\eqref{an04ce5v1:envelope} then gives $\|U_2\|+\|W_2\|+\|b_2\|\le r(t)=C_0q e^{Q(t-t_0)/2}$ up to $t_R$, with $r(t_R)\le C_1e^{-LR}$ independently of fixed $R$ for small $q$.

\section{Uniform prefix estimates}\label{an04ce5v1:prefixsec}
The following details make the comparisons quantitative before the
terminal clock. Set $d_i=W_i-b_i$ and $w=Q(1-n)$. In the characteristic
equations $U'=\mathsf R^\top W,\ W'=\mathsf R U$, the regression
decomposition, including its uniform per-receiver shape error, gives
\[
 \mathsf R=\tfrac12 Q_1(1-m,-z)^\top v_1^\top+
            \tfrac12 Q(-h,1-n)^\top v_2^\top+O(q).
\]
Here the pairs denote column vectors in the outer products. The error
follows by data Cauchy--Schwarz against $D$ times the receiver gate and
the $O(q(1+S))$ cluster shape residual. Equivalently the output-layer
equation uses $b_i$ with error $O(q^2|U|)$ by orthogonality also of the
feature remainder. We will often only need $O(q|U|)$.
In particular $\mathsf R_{22}\ge-Cq$ and
$|\mathsf R_{12}|+|\mathsf R_{21}|\le C(|z|+|h|+q)$ uniformly.

For $x=1-m$, the scalar damped equations from the four-regression flow are
\[
 x'=-a_m x+O(|T_1 z|+|h|C+q),\qquad
 z'=-a_z z+\tfrac12 w C+\tfrac12 Q_1T_1 x+O(q).
\]
Both coefficients have positive upper and lower bounds. On the incoming
interval $a_m\ge Q_1m^2/S>L+\eps$ for a fixed $\eps>0$.
The integrated cost of $T_1$ times the homogeneous entrance deficit
is $O_{b_0}(q)$: $|T_1|\le Cr(t)$ on the incoming interval, and the final
fixed interval starts after the entrance deficit has decayed to $o(q)$.
The coefficient comparison can use one \emph{fixed} incoming cutoff
before $R$ is chosen. Integrating the $z$ propagator up to an arbitrary
$t$ gives weight comparable to $1\wedge(t-u)$. Duhamel in the $x$
equation bounds the remaining $T_1x$ cost by
$O(B(t))+o_q(1)(I(t)+\int_{t_0}^t |z|)$. Thus
\begin{align}
 \int_{t_0}^t |z|&\le C(I+B),&
 \int_{t_0}^t z_-&\le o_q(1) I+CB, \nonumber\\
 \int_{t_0}^t z&\ge c\int_{t_0}^t[1\wedge(t-u)]C(u)\,du-o_q(1)I-CB.
 \label{an04ce5v1:prefix}
\end{align}
Constants here and in the next estimates are independent of the compact
right-label truncation and the large incoming offset $R$.

Differentiating $b_i$ uses $v_i$, with $|v_1^\top v_2|\le Cq$,
$b_i=(U_i)_++O(q|U|)$ and
$|b_i(1-|v_i|^2)|\le Cq|U|$. The last bound uses
$\sup_{x>0}x\Phi(-x)<\infty$ in the positive coordinate.
The $d_1$ equation contains the terms
\[
 2d_1'=-Q_1x|v_1|^2d_1-Qb_2\langle b_2,d_1\rangle+
 Q_1 z W_2 |v_1|^2-Qb_2\langle b_1,b_2\rangle+O_{L^2}(q).
\]
At $t_0$, $\norm{d_1}=O_{b_0}(q)$: integrate the same inequality from
initialization, with $\norm{d_1(0)}=O(\sqrt{S_0})$, the HBD bounds for
the cross regressions and overlap $O(H_0V_0+qH_0^2)$, and the integrated
early shape cost $O(q\int H_0)=O_{b_0}(q)$, with the entrance envelopes of Section~\ref{an04ce5v1:setup}. The growth from $x_-$ is bounded by the early mass integral.
Dropping the dissipative rank-one term and using $\int |x|\le C$ gives
\[
 \norm{d_1(t)}\le C(I+B),\qquad
 \norm{d_1(t)}\le C(r(t)I+B)\quad(t\le t_R).
\]
Since $h=\langle b_1,b_2\rangle+\langle d_1,b_2\rangle$,
$\langle b_1,b_2\rangle=C+O(q)$ and $\int\norm{b_2}\le C$, we get
\begin{equation}
 \int_{t_0}^t|h|\le C(I+B),\qquad
 h_-(t)\le C r(t)^2 I(t)+C B(t)\quad(t\le t_R).
 \label{an04ce5v1:earlyh}
\end{equation}
Similarly $2d_2'=-w|v_2|^2d_2-Q_1zb_1+QhW_1|v_2|^2+O_{L^2}(q)$, so
$\norm{d_2}\le C(I+B)$. This dissipative bound will also control costly
returns outside the retained bulk.

For clarity, the terminal weight in \eqref{an04ce5v1:prefix} can be removed
up to constants and $O(B)$. Differentiation of
$C=\int(U_1)_+(U_2)_+\dd$ a.e. gives $|C'|\le C_2(C+I+B)$:
replace $W_i$ in each diagonal characteristic term by $b_i+d_i$,
use the two norm bounds, and use
$|z|+|h|\le C_2(C+I+B)$ for the off-diagonals.
With $I'=C$, Gronwall on an interval of length at most two,
starting at a point of $[t-2,t-1]$ where $C\le I(t-1)$, gives
$I(t)+C(t)\le C_3(I(t-1)+B(t))$ for $s\ge2$. Hence the weighted integral
is at least $cI(t)-CB(t)$.

\paragraph{Dominance over the accumulated Gaussian errors.}
We need $B/I$ small uniformly, even at delays growing slower than $\ell$.
In precise terms
\begin{equation}
 \lim_{b\to\infty}\limsup_{q\to0}\ \sup_{t_h+b\le t\le T} B(t)/I(t)=0 .
 \label{an04ce5v1:uniform}
\end{equation}
Here are details of the moving-band check. On a contrary sequence with
$s=t-t_0\to\infty$, $t\le t_R$, one would have $I(t)=O(q(1+s))$,
so $\int_{t_0}^t(|z|+|h|+q)=O(q(1+s))$. Choose the initialized Q1 band
$\eps\le\cos\alpha\le2\eps$, $\eps=(1+s)^{-4}$.
At $t_0$ both horizontal components are positive and comparable to $\eps$,
and both vertical components are positive and comparable to $q$.

To verify this moving entrance, let
$H_0=\sqrt{S_0}e^{(1/2+q^2)t'}$ and
$V_0=\sqrt{S_0}e^{(L+q^2)t'}$ for $0\le t'\le t_0$.
HBD gives $|U_2|+|W_2|\le C_{b_0}V_0$ and off-diagonal field
$O_{b_0}(q+H_0V_0)$. The integrated perturbation in horizontal units is
$O(\int (V_0/H_0)(q+H_0V_0))=O(q)$; hence horizontal size is
$O(\eps H_0)$. The vertical perturbation in units of $V_0$ is then
$O(\eps\int (H_0/V_0)(q+H_0V_0))=O(\eps)$.
For lower comparison the unperturbed symmetric pairs start with equal
positive components; the integrated fitted diagonal is $O(\int S)=O(1)$.
On the matching lower faces, strong standardized gate margin is
$\gtrsim\eps/q$ and weak margin is
$\gtrsim\min(1/q,e^{(1-\lambda)(t_0-t')/2}/\eps)$.
The integrated weak wrong-gate cost is bounded by
$C\ell e^{-c_1/q^2}+C e^{-c_1/\eps^2}$, including very slowly growing
$s$. These bounds close the lower faces by variation of constants for the
symmetric and antisymmetric combinations; the additive perturbations are
$o(\eps)$ and $o(1)$ in the respective units.

After $t_0$ the horizontal diagonal has bounded absolute integral.
Until $t$, the vertical growing factor is comparable to $e^{Q(t'-t_0)/2}$
as $\int|n|=O(1)$, the weak standardized margin on the trial faces is
$\gtrsim1/\eps$, and $q s=o(1)$. The horizontal perturbation on these
faces is $O(q(1+s))=o(\eps)$ since $qe^{Qs/2}=O(1)$. The discounted
vertical perturbation is $O(\eps q(1+s))=o(q)$. Initial imbalance of
the two horizontal components is $o(\eps)$ by balance and their positive
entrance, and the vertical growing pair has a positive lower comparison.
Thus the band keeps horizontal components $\asymp\eps$ and vertical
components $\asymp q e^{Q(t'-t_0)/2}$. On the final unit its overlap alone
integrates to $\gtrsim \eps^2 q e^{Ls}$, contradicting the assumed bound.
For $t\ge t_R$ use monotonicity of $I$. The $q^{3/4}$ lower bound
at $t_R-1$ can also be checked with the same proof: if
$I(t_R-1)=O(q^{3/4})$, use $\eps=q^{1/3}$ instead. The horizontal
perturbation is $O(q^{3/4})=o(\eps)$ and the discounted vertical one
is $O(q^{13/12})=o(q)$. On the last unit the vertical amplitude is
bounded below by a positive constant depending on $R$. The band alone
then integrates to $\gtrsim_R q^{2/3}$, a contradiction.
This proves \eqref{an04ce5v1:uniform}. Together with
\eqref{an04ce5v1:prefix} it gives, uniformly after a sufficiently late fixed
delay, $\int z\ge c I$ and $\int z_-/I$ as small as desired.

\section{Choosing a bulk without losing exposure to discarded labels}\label{an04ce5v1:retained}
In the following estimates the comparison constants for a retained label
are relative to $k_\alpha=\cos\alpha>0$. They can be kept independent of
how close the compact truncation is to the quadrant boundaries. This
point will allow the discarded mass to be chosen last.

In $A_R$, remove initial strips of total length $\delta$ at the axes.
For the remaining compact Q1 and Q4 arcs, the fixed-delay HBD/BRW
horizontal entrance is asymptotic to $2\sqrt{M(b_0)}\,k_\alpha$ on both
layers. Errors can be taken small relative to $k_\alpha$ on each such
compact. On $[t_0,T]$ the integral of the supremum of the two
off-diagonal field entries is $o_q(1)$. Indeed the incoming envelope
bounds the integral of $|z|+|h|$ before $t_R$ by $Ce^{-LR}$, and on the
remaining fixed interval they tend uniformly to zero. Take large fixed
$R$ for this assertion only. The horizontal pair has diagonal coupling
$\mathsf R_{11}$ of bounded absolute integral. Its symmetric mode keeps
a fixed multiple of its initial size, and its antisymmetric entrance tends
to zero. Its integrated perturbation is $o_q(1)$ in the label norm of
the time supremum by the bound on $\mathcal U$.

Also $\norm{\mathbf1_{A_R}\sup_{[t_0,T]}(|U_2|+|W_2|)}=o_q(1)$.
To pass from the supremum of norms to this stronger norm, integrate
the vertical equations: the off-diagonal cost just bounded is $o_q(1)$;
the integral of the right vertical norm is $o_q(1)$ by the incoming
envelope up to $t_R$, then uniform axis separation on the last bounded
interval, then $R\to\infty$. Excluding an additional label set of measure
$o_q(1)$ by Chebyshev on each fixed compact therefore leaves a set $K$
with, for an absolute comparison constant $D$ (fixed family),
\begin{equation}
 D^{-1}k_\alpha\le U_1,W_1,b_1\le D k_\alpha,\qquad
 |U|+|W|\le D k_\alpha\qquad(\alpha\in K,\ t_0\le t\le T).
 \label{an04ce5v1:bulk}
\end{equation}
For example, remove errors bigger than a small fixed multiple of the
minimum $k_\alpha$ on the compact. Moreover,
$\norm{\mathbf1_{A_R\setminus K}\mathcal U}\le C\sqrt\delta+o_q(1)$:
the horizontal entrance is uniformly bounded even on the initial strips,
its later amplification is bounded, and the integrated horizontal
perturbation and right vertical supremum have vanishing norm.
These estimates allow the set $K$ to depend on the full path. All
barriers below concern this fixed set of labels for the given $q$.

For $A_L$, the initial horizontal norm at $t_0$ is $O_{b_0}(q\ell)$
(HBD's left-ancestry estimate). Horizontal comparison on any prefix gives
\begin{equation}
 \norm{\mathbf1_{A_L}\sup_{[t_0,t]}(|U_1|+|W_1|)}
 \le C(q\ell+I(t)).
 \label{an04ce5v1:left}
\end{equation}
It is also $o_q(1)$ uniformly, by the vanishing off-diagonal integral.
The left vertical difference $W_2-U_2$ has vanishing norm of the time
supremum. Its coefficient is $-\mathsf R_{22}\le Cq$, its initial norm
is $O_{b_0}(q)$, and the forcing is bounded by the off-diagonals times
the horizontal components. Negative $U_2$ has vanishing norm uniformly
by the same negative-excursion estimate as at $T$. In particular
$\norm{\mathbf1_{A_L} d_2}=o_q(1)$ uniformly.

One further consequence will pay for returns of weak labels. On
$D_L=A_L\cap\{p>0\}$, the terminal upper face gives
$U_2(T)\le Cq|U(T)|$, so $\norm{\mathbf1_{D_L}U_2(T)}=o_q(1)$ using
the negative-excursion bound as well. Trace $(U_2+W_2)/2$ backwards.
The inverse propagator for its coefficient $\mathsf R_{22}$ is bounded
by $e^{Cq(T-t_0)}$ and the integrated off-diagonal forcing has vanishing
norm. Combining with the difference and horizontal bounds proves
\begin{equation}
 \norm{\mathbf1_{D_L}\mathcal U}=o_q(1).
 \label{an04ce5v1:activeleft}
\end{equation}

\section{Entry into the negative-input region}\label{an04ce5v1:entry}
Write $\nu=\inf_{[t_0,T]}(1-n)>0$ with a common lower bound for small
$q$. For $\alpha\in K$, \eqref{an04ce5v1:bulk} and the prefix estimates give,
after a late fixed delay $b$,
\[
 \int_{t_0}^t z b_1\,du\ge c_1 k_\alpha I(t)
 \qquad(t_h+b\le t\le T).
\]
Indeed the positive part is multiplied by at least $k_\alpha/D$ and
the negative part by at most $D k_\alpha$. The constant $c_1$ works
independently of $\delta$. We next pay for weak excitation up to
$t_b=t_h+b$. At each fixed $b$ the scaled weak-feature integral on a
retained right compact satisfies
\[
 \limsup_{q\downarrow0}\sup_K
 \frac{1}{q k_\alpha}\int_{t_0}^{t_b}|b_2|\,du
 \le C_\delta(1+b)e^{(L-\nu_\delta)b},\qquad \nu_\delta>0.
\]
This follows by $|b_2|\le (U_2)_++Cq|U|$, fixed-interval tagged
transfer, and the COD rank ceiling for $x_+$ on ranks bounded away from
the upper endpoint; the Q4 bottom path obeys the same bound. Take
$\nu_\delta<L$ if necessary. In contrast BRW's protected positive
first moment on the last unit gives
$\liminf_{q\to0} I(t_b)/q\ge c e^{Lb}/(1+b)$ for all large fixed $b$.
To transfer this lower bound, retain the positive overlap from the
right labels and use fixed-time extraction:
$(U_1)_+(U_2)_+/q$ converges there to the radial mass times $x_+$;
the nonnegative remainder may be discarded. TBW's bulk upper input
face also puts all retained compacts strictly below $U_2=0$ at $t_b$.
Thus, after fixing the compacts, choose $b$ sufficiently large for
the initial weak integral and $|W_2(t_0)|$ to cost less than any
prescribed fixed fraction of $k_\alpha I(t_b)$.

As long as $U_2\le0$ after this entrance on a retained label,
$|b_2|\le Cq k_\alpha$. The output equation integrates to
\begin{equation}
 W_2(t)\le -c_2 k_\alpha I(t),\qquad t_b\le t\le T,
 \label{an04ce5v1:wnegative}
\end{equation}
on this stopped interval: the additional weak cost is $Cq(t-t_b)k_\alpha$,
absorbed uniformly by \eqref{an04ce5v1:uniform}; the output shape error is
smaller. The constant $c_2>0$ is independent of $\delta,R,b$ once
the entry requirements are met. The input equation on that interval is
\begin{equation}
 U_2'=\tfrac Q2[(1-n)W_2-hW_1]\theta_2+O(qk_\alpha).
 \label{an04ce5v1:gate}
\end{equation}
Fix $\kappa>0$ sufficiently small compared with $\nu c_2/D$.
Whenever $h_-\le\kappa I$, \eqref{an04ce5v1:gate} gives
$U_2'\le-c_3 k_\alpha I\theta_2+Cq k_\alpha$.

Up to $t_R$ this bound on $h_-$ follows from \eqref{an04ce5v1:earlyh}:
choose $R$ large first, then make $B/I$ small by the fixed delay $b$.
At $U_2=0$, $\theta_2=1/2$. The strict inward sign protects the face
for every retained label up to $t_R$. In fact, by $t_R$ the stronger
face
\begin{equation}
 U_2\le -d q k_\alpha\sqrt\ell
 \label{an04ce5v1:deepgate}
\end{equation}
holds, with a positive $d$ independent of $\delta$.
On $[t_R-1,t_R]$, $I\gg q^{3/4}$. In the strip between
\eqref{an04ce5v1:deepgate} and zero, $\theta_2\ge c q^\gamma$ for some
$\gamma<1/4$ by choosing $d$ small in the Gaussian lower-tail bound.
Here $|U|\ge k_\alpha/D$; restricting to $D_2>\rho/2$ costs only an
exponentially small probability in the weighted gate. The downward
speed is at least $c k_\alpha q^{3/4+\gamma}$, larger than the strip
width per unit time. The face is reached in $o(1)$ time and is inward
up to $t_R$.

\section{Closing the final interval}\label{an04ce5v1:centered}
Bootstrap \eqref{an04ce5v1:deepgate} for the retained right bulk. On that
face and below, $\theta_2=o_q(1)$ on $K$ by the Gaussian upper tail
and $|U|\le Dk_\alpha$. We use the following signed
regression approximation on the stopped interval:
\begin{equation}
 |T_2-h|\le (C\sqrt\delta+o_q(1))I(t),\qquad t_R\le t\le T,
 \qquad T_2=\int W_1 W_2 |v_2|^2\dd .
 \label{an04ce5v1:t2}
\end{equation}
In fact $T_2-h=\int W_1|v_2|^2 d_2\dd+O(q)$ since
$|b_2(1-|v_2|^2)|\le Cq|U|$. On $K$, $|v_2|^2=o_q(1)$.
On the other right labels $\norm{W_1}\le C\sqrt\delta+o_q(1)$.
Use $\norm{d_2}\le C(I+B)$ on both sets. On the left, use
\eqref{an04ce5v1:left} for $W_1$ and the vanishing left norm of $d_2$.
Since $I\gg q^{3/4}$, all additive $O(q\ell)$ errors are negligible.
This argument needs small right \emph{physical weight} in the exception.

The exact scalar regression equation now reads
\[
 h'=\tfrac Q2(1-n-\norm{\nabla h}^2)h
       +\tfrac w2(T_2-h)+\tfrac{Q_1}2(1-m)C+O(q).
\]
On this fixed centered interval $1-m=o_q(1)$ and $C\le C_3(I+B)$
by the two-unit estimate above. The homogeneous coefficient is at most
$Q/2$, since $n>0$ uniformly here for each fixed $R$ and small $q$.
The entry \eqref{an04ce5v1:earlyh} gives
$h_-(t_R)\le(Ce^{-2LR}+o_q(1)) I(t_R)$. As
$T-t_R=R+O(1)$, the initial negative part after amplification is at
most $(C'e^{-LR}+o_q(1))I(t)$, with $C'$ independent of $R$.
Integration of \eqref{an04ce5v1:t2} adds $(C_R\sqrt\delta+o_q(1))I(t)$,
using monotonicity of $I$. Therefore choose $R$ large enough that
the first coefficient is below $\kappa/4$, then choose $\delta$ small
enough for the second to be below $\kappa/4$. The entry delay $b$ is
chosen \emph{after} this truncation, and $q$ last.
This proves $h_-<\kappa I$ with slack. Equations
\eqref{an04ce5v1:wnegative} and \eqref{an04ce5v1:gate} now point strictly inward on
\eqref{an04ce5v1:deepgate}, with the same $q^\gamma$ lower bound there.
The estimates hold uniformly also on a thin outer strip about the face
(reduce $d$ slightly for the upper-tail test), so first exit is excluded
even if a supremum over the retained labels is approached rather than
attained. This closes the bootstrap through $T$.

\section{Subtracting all positive output outside the core}\label{an04ce5v1:endpoint}
On $K$, terminal exposure is at least $c q k_\alpha\sqrt\ell$, and
$-W_2(T)\ge c_2 k_\alpha I(T)$. The retained Q4 contains, for small
$\delta$ and small $q$, all but vanishing measure of a fixed central
Q4 subarc. Thus its favorable contribution is at least
$c_4 q\sqrt\ell\,I(T)$ with $c_4>0$ independent of the truncation.
Every retained Q1 label also has negative output.

To charge \emph{all} adverse labels, including any exceptional Q4,
let $D_*=\{p>0\}\setminus K$. Equations \eqref{an04ce5v1:activeleft} and
the right-strip estimate give
$\norm{\mathbf1_{D_*}\mathcal U}\le C\sqrt\delta+o_q(1)$.
The exposure excursion bound consequently gives
$\norm{\mathbf1_{D_*}p}\le(C\sqrt\delta+o_q(1))q\sqrt\ell$.
On the probe-active side $b_2\le Cq|U|$ because $U_2\le Cq|U|$.
Hence $\norm{\mathbf1_{\{p>0\}}(W_2)_+}\le\norm{d_2}+Cq\sqrt S
\le C I(T)$. Cauchy--Schwarz pays the entire adverse output by
$(C_5\sqrt\delta+o_q(1))q\sqrt\ell\,I(T)$.
Take $\delta$ small enough also that $C_5\sqrt\delta<c_4/2$.
The constants here precede the compact-specific choice of $b$.

Finally let $G_\alpha=\int_0^T[Q_1zb_1-w b_2]\dt$ and define
$M_q=(\int_{\rm Q4}pG_\alpha\dd-\int_{{\rm Q4}^c}p(-G_\alpha)_+\dd)/(2q^2)$.
The output projection error in replacing $-W_2(T)$ by $G_\alpha/2$
in these weighted integrals is $o(q^2)$ in absolute integral. Indeed
the shape error costs $O(q^2\ell)\norm p$, and the initial term costs
$O(\sqrt{S_0})\norm p$. Positive part is 1-Lipschitz. Consequently
\[
 M_q\ge c\sqrt\ell\, I(T)/q-o(1),\qquad
 H^2=O(\ell),\qquad I(T)/q^{3/4}\longrightarrow\infty .
\]
Since $\sqrt\ell\,I(T)/q\gg\ell$, this proves \eqref{an04ce5v1:criterion} with diverging margin. Directly, if
$F_i=f_{T,i}(\xi)$, then
$e_q=H^2/2-q^{-1}\sin(q\zeta)H+
(\cos(q\zeta)+F_2/2)F_2/q^2$.
The last coefficient tends to one by $\norm p=O(q\sqrt\ell)$.
Thus $e_q(T)\to-\infty$, which gives the strict comparison with
the finite coherent minimum at the measured learned endpoint. Since $F_i=O(q\sqrt\ell)$ the unscaled probe loss tends to $1/2$, as asserted in Theorem~\ref{an04ce5v1:main}.
\end{document}
